
\subsubsection{2.1 SGD Learning Dynamics}

The Frequency Principle \citep{Xuetal2019} — also termed spectral bias \citep{Rahamanetal2019} — establishes empirically and theoretically that neural networks trained with gradient descent learn lower-frequency components of target functions before higher-frequency ones. This result has been replicated across architectures and attributed to properties of the neural tangent kernel. Spatial frequency in image features directly maps to a notion of complexity: large-scale textures (background patterns) are lower frequency and simpler; fine-grained structural features (morphological detail) are higher frequency.

Shah et al. [2020] characterize the Simplicity Bias: with high probability, SGD finds the simplest classifier consistent with training data, preferentially encoding linearly separable features over features requiring nonlinear separation. Together, the Frequency Principle and Simplicity Bias predict that if spurious features are simpler than core features, they will be encoded earlier during training — the central prediction tested in this work.

Both results are established on synthetic or controlled tasks. Neither paper provides a measurement protocol for the temporal dynamics of spurious versus core feature learning on standard spurious correlation benchmarks such as Waterbirds or CelebA. The present work operationalizes these theoretical predictions in that setting.

\subsubsection{2.2 Spurious Correlation Benchmarks and Mitigation Methods}

Sagawa et al. [2020] formalize spurious correlations as a group robustness problem and introduce GroupDRO, which minimizes worst-group training loss using group annotations. Liu et al. [2021] propose JTT, which identifies minority samples by their misclassification under a first-pass ERM model and upweights them in a second training run. Kirichenko et al. [2022] show that DFR — refitting only the last linear layer with class-balanced data — achieves state-of-the-art WGA without group annotations at training time. The stated rationale for DFR is that ERM backbones already encode core features; DFR recovers them by rebalancing the final classifier.

None of these methods quantify when during training spurious or core features are encoded, or characterize the gradient mechanism driving differential learning speed. The present framework provides this temporal characterization.

\subsubsection{2.3 Early Training Dynamics and Shortcut Learning}

Mangalam and Girshick [2021] observe that shortcut features emerge in early training phases, but this observation is qualitative and lacks a systematic measurement protocol, statistical validation across seeds, or gradient-level mechanistic characterization. Frankle and Carlin [2019] identify early-epoch decisions as critical via the lottery ticket hypothesis. Jiang et al. [2020] characterize easy and hard examples through per-sample training dynamics. Toneva et al. [2019] study forgetting events during training. These works collectively establish that early training dynamics encode meaningful information about spurious feature acquisition, but none defines a reproducible protocol for measuring the temporal competition between feature types.

\subsubsection{2.4 Linear Probing as a Measurement Tool}

Linear probing — training a linear classifier on frozen representations — has been applied to track feature emergence in language models \citep{Tenneyetal2019}, evaluate representation quality in self-supervised learning \citep{Chenetal2020}, and understand intermediate layer representations [Alain and Bengio, 2016]. These applications treat probing as a static evaluation tool. Applying checkpoint linear probing at regular intervals throughout supervised training to obtain a time series of feature-type-specific probe accuracies — the core measurement innovation here — has not been previously reported in the context of spurious correlation benchmarks.

\subsubsection{2.5 Positioning}

\begin{table}[htbp]
\centering
\begin{tabular}{lll}
\toprule
Prior Work & Contribution & Limitation Relative to This Work \\
\midrule
Xu et al. [2019], Rahaman et al. [2019] & Frequency Principle & Synthetic tasks; no spurious correlation benchmarks \\
Shah et al. [2020] & Simplicity Bias & Synthetic settings; no measurement protocol \\
Sagawa et al. [2020] & GroupDRO, Waterbirds/CelebA & Focuses on intervention, not temporal measurement \\
Liu et al. [2021] & JTT & No measurement of when shortcuts form \\
Kirichenko et al. [2022] & DFR & No temporal analysis of core feature emergence \\
Mangalam \& Girshick [2021] & Shortcuts emerge early (qualitative) & No systematic protocol; no statistical validation \\
\bottomrule
\end{tabular}
\end{table}

